% R008 — Guide to Cultivating Complex Analysis (Bahasa Indonesia)
% Reader-facing translation unit: source/ca-v1.9/ca.tex lines 1400–1412.
% Source release: v1.9 (commit a4d25d9ba37a486a5a020219eb8bd4e6602fd14c).
% Translation provenance: OpenAI Codex gpt-5.6-sol, Ultra, acting on the user's request.
% Preserve source equations, notation, labels, references, and comments.

\begin{exbox}
\begin{exercise} \label{exercise:exponetoonestrip}
Buktikan bahwa identitas-identitas dalam \eqref{eq:exponentialidentities} adalah satu-satunya
identitas semacam itu dengan menunjukkan bahwa jika $2k\pi < \Im z \leq 2(k+1)\pi$ dan
$2k\pi < \Im w \leq 2(k+1)\pi$, maka $e^z=e^w$ mengakibatkan $z=w$.
\end{exercise}

\begin{exercise}%[\neededexmark]
Gunakan $e^{z+w} = e^z e^w$ dan $e^0 = 1 \neq 0$ untuk menunjukkan bahwa $e^z \neq 0$ untuk setiap
$z \in \C$. Dengan kata lain, tunjukkan bahwa jika suatu fungsi $f$ memenuhi $f(z+w)=f(z)f(w)$
dan $f(0) = 1$, maka $f(z) \neq 0$ untuk setiap $z$.
\end{exercise}
\end{exbox}
